\begin{claim}{C.14}{Proved under stated hypotheses}{Pure-injective resolutions close a further part of the raw-adjoint gap}\label{C14}
Assume the normalized hypotheses of Claim~\ref{C6}, including the interval condition (d). Suppose that every map
\[
 F(\alpha\otimes X):F(X)\longrightarrow F(J\otimes X)
\]
is degreewise pure as a map of complexes of $k$-modules. Suppose also that some integer $d\geq0$ bounds the pure-injective dimension of every projective $k$-module. Then the underived adjoints $F,G$ induce quasi-inverse equivalences, as in Claim~\ref{C7}.

The purity hypothesis holds if the tensor product on $\mathcal M$ is compatible with scalar copowers: the canonical comparisons
\[
 (U\otimes_k\mathbf1)\otimes X\ \cong\ U\otimes_kX
 \qquad(U\in\Mod k,\ X\in\mathcal M)
 \tag{C-copower}\label{Ccopower}
\]
are isomorphisms, compatibly with the unit, symmetry, and associativity. In particular this compatibility holds when the internal tensor product preserves colimits in each variable. It is an additional hypothesis here: mere $k$-bilinearity does not assert preservation of cokernels or infinite coproducts.

With this compatibility, the conclusion applies to every ring of finite pure global dimension. It therefore applies to every commutative ring of cardinality at most $\aleph_t$ for a finite integer $t$, in particular every countable commutative ring, by the Gruson--Jensen cardinal bound stated in \citet[Theorem 2.6.3]{BensonGnacadja1999}. No cardinal restriction on the modules, the generator category, or the resolving functor is needed.
\end{claim}
\begin{proof}
We recall the exact meaning of the extra module-theoretic assumption. A monomorphism $A\to B$ of $k$-modules is pure if $U\otimes_k A\to U\otimes_k B$ is monic for every $k$-module $U$. A module $I$ is pure-injective if every map $A\to I$ extends across every such monomorphism. Its pure-injective dimension is at most $d$ if it admits a pure-exact resolution
\[
 0\longrightarrow P\longrightarrow I^0\longrightarrow\cdots
 \longrightarrow I^d\longrightarrow0
\]
with all $I^j$ pure-injective. These resolutions exist without the length bound; for example the character double-dual map is a pure monomorphism into a pure-injective module, and we iterate on its cokernel. The extension property, double-dual construction, and equivalence of the global pure-projective and pure-injective bounds are recorded in \citet[Sections 2.3--2.5]{BensonGnacadja1999}.

Let $A=\prod_{n\geq1}D^n(I_n)$, where every $I_n$ is pure-injective. A chain map from a complex $C$ to $D^n(I_n)$ is specified by a module map $C_{n-1}\to I_n$. Consequently a map $C\to A$ extends across every degreewise pure monomorphism $C\to C'$: extend these component maps and combine them by the universal property of the product. In particular the counit
\[
 \varepsilon_A:FGA\longrightarrow A
\]
extends across the degreewise pure map $F(\alpha\otimes GA)$ to a map $F(J\otimes GA)\to A$. Under $F\dashv G$, this extension corresponds to a map $r:J\otimes GA\to GA$ whose composite with $\alpha\otimes GA$ is $1_{GA}$. This assertion follows from naturality of the adjunction and the fact that the counit corresponds to $1_{GA}$. Thus $GA$ is a retract of $J\otimes GA$, so $FGA$ is acyclic by (d). This proves the required vanishing for products of disks on pure-injective coefficients; it does not call a pure monomorphism split.

Now let $P$ be an acyclic nonnegative complex of projective $k$-modules. Claim~\ref{C1} gives $P\cong\prod_{n\geq1}D^n(P_n)$ with each $P_n$ projective. Choose a pure-injective resolution of length at most $d$ of each $P_n$, padding by zero if necessary. Applying $D^n$ and taking the product gives an exact sequence of complexes
\[
 0\longrightarrow P\longrightarrow A^0\longrightarrow\cdots
 \longrightarrow A^d\longrightarrow0,
 \qquad A^j=\prod_{n\geq1}D^n(I_n^j).
\]
There are only two disk contributions in each degree, so exactness here does not require an infinite-product exactness argument. Both $F$ and $G$ are exact. Hence applying $FG$ gives a finite exact sequence, all of whose terms $FGA^j$ are acyclic by the preceding paragraph. The homology long exact sequences for its successive kernels give $FGP$ acyclic. This is precisely the remaining test of Claim~\ref{C10}; normalization already proves the counit test on projective modules in degree zero. Claim~\ref{C3} now proves the asserted equivalence by the raw adjoints.

We check the sufficient tensor hypothesis rather than identify it with $k$-bilinearity. Every $F_n$ is cocontinuous and $k$-linear, so it commutes with scalar copowers. Under \eqref{Ccopower}, symmetry and associativity identify $U\otimes_k(\alpha\otimes X)$ with $\alpha\otimes((U\otimes_k\mathbf1)\otimes X)$. The latter is monic by (d). Applying the exact functor $F_n$ identifies its image with $U\otimes_kF_n(\alpha\otimes X)$, proving purity for every $U$. If the internal tensor is cocontinuous, \eqref{Ccopower} itself follows by writing $U$ as the cokernel of a map between free $k$-modules: both sides preserve the resulting coproducts and cokernel, and the comparison is the unit isomorphism on each free summand.

Finally finite pure global dimension gives the required uniform pure-injective bound by its equivalent injective definition. The cited cardinal theorem supplies such a bound $d=t+1$ when $|k|\leq\aleph_t$; finite rings are included in the countable case. This is a direct application of that theorem to coefficient modules, not a bound on the cardinalities of our diagrams. The zero ring gives only zero modules and the conclusion is immediate there.
\end{proof}

